\section{总结与延伸}
\subsection{核心总结表}
\begin{tabularx}{\textwidth}{lXl}
\toprule
维度 & 核心洞察 & 实践启示 \\
\midrule
奖励信号 & 图像、语言、success classifier 取代人工 reward & 设计 reward pipeline 时加入置信度/ensemble 机制； \\
环境运维 & Forward-backward controllers 与 time-aware buffer 实现 reset-free & 每个 reset policy 都需要 recovery 记录并与 scheduler 绑定； \\
目标管理 & Goal-Query 与 curriculum policy 实现难度自适应 & 维护 goal repository，并记录 novelty/skill dependency； \\
安全与运维 & safety penalty + regular monitoring 防止 drift & 搭建 dashboard，保留 rollback 通道； \\
\bottomrule
\end{tabularx}

\subsection{进一步阅读}
\begin{itemize}
\item Sharma et al., ``Autonomous Reinforcement Learning: Formalism and Benchmarking,'' ICLR 2022
\item Eysenbach et al., ``Leave No Trace: Learning to Reset for Safe and Autonomous RL,'' ICLR 2018
\item Pong et al., ``Skew-Fit: State-Covering Self-Supervised RL,'' ICML 2020
\item Nair et al., ``Visual Reinforcement Learning with Imagined Goals,'' NeurIPS 2018
\item Ma et al., ``VIP: Towards Universal Visual Reward and Representation via Value-Implicit Pre-Training,'' ICLR 2023
\end{itemize}

\subsection{本章小结}
通过奖赏、reset、goal、安全与监控五个维度的系统讲解，本讲展示了自主机器人学习的完整路径，并为部署级 Autonomy 提供了可执行的反馈/回退机制。

\end{document}
